% ==============================================================
\section{問題定義}\label{sec:problem}
% ==============================================================

\subsection{記法}

$\mathcal{I} = \{i_1, i_2, \ldots, i_m\}$ を有限のアイテム全集合，
$D = [d_1, d_2, \ldots, d_N]$ を時刻順のトランザクション列とする．
各トランザクション $d_t$（$t = 1, \ldots, N$）は $\mathcal{I}$ の部分集合である．
アイテムセット（パターン）$P \subseteq \mathcal{I}$ が
トランザクション $d_t$ に\emph{出現}するとは $P \subseteq d_t$ が成立することをいう．

\begin{definition}[サポート時系列]\label{def:support_series}
パターン $P$，ウィンドウサイズ $W \in \mathbb{Z}_{>0}$，
トランザクション列 $D$（長さ $N$）に対して，
$P$ の\emph{サポート時系列} $s_P = [s_P(0), s_P(1), \ldots, s_P(N-W)]$ を
\begin{equation}
s_P(t) = \bigl|\{d_j \in D \mid t \leq j < t + W,\; P \subseteq d_j\}\bigr|
\label{eq:support_series}
\end{equation}
と定義する．これは位置 $t$ を始点とするウィンドウ $[t, t+W)$ 内で
$P$ を含むトランザクション数である．
\end{definition}

\begin{definition}[外部イベント]\label{def:event}
\emph{外部イベント} $e = (\textit{id}, \textit{name}, t_s, t_e)$ は，
一意な識別子，名称，およびイベントが有効な期間を表すタイムスタンプ区間 $[t_s, t_e]$ から成る．
既知の外部イベント集合を $\mathcal{E} = \{e_1, e_2, \ldots, e_K\}$ とする．
\end{definition}

\begin{definition}[変化点（閾値交差）]\label{def:change_point}
本稿では，サポート時系列 $s_P$ に対する\emph{変化点}として，
一般のレベルシフト点すべてではなく，
閾値 $\theta$ に関する密状態／非密状態の切替点のみを扱う．
位置 $\tau \in \{0, 1, \ldots, N-W\}$ が変化点であるとは，
次のいずれかが成り立つことをいう:
\begin{align}
s_P(\tau) \geq \theta \land (\tau = 0 \lor s_P(\tau-1) < \theta), \label{eq:up_crossing}\\
s_P(\tau) < \theta \land \tau > 0 \land s_P(\tau-1) \geq \theta. \label{eq:down_crossing}
\end{align}
式~(\ref{eq:up_crossing}) を満たすとき方向を「上昇」，
式~(\ref{eq:down_crossing}) を満たすとき方向を「下降」とする．

さらに，$T = N-W+1$ とし，
\begin{align*}
\mathcal{L}^{-}_{\tau} &= \{\max(0, \tau-h), \ldots, \tau-1\},\\
\mathcal{L}^{+}_{\tau} &= \{\tau, \ldots, \min(T-1, \tau+h-1)\}
\end{align*}
を交差前後の参照窓とする．
変化量はこれらの局所平均差
\begin{align*}
\bar{s}^{-}(\tau) &=
\frac{1}{\max(1, |\mathcal{L}^{-}_{\tau}|)}
\sum_{j \in \mathcal{L}^{-}_{\tau}} s_P(j), \\
\bar{s}^{+}(\tau) &=
\frac{1}{\max(1, |\mathcal{L}^{+}_{\tau}|)}
\sum_{j \in \mathcal{L}^{+}_{\tau}} s_P(j)
\end{align*}
に基づき
\begin{equation}
\text{magnitude}(\tau) = |\bar{s}^{+}(\tau) - \bar{s}^{-}(\tau)|
\end{equation}
と定義する．
したがって，密集区間がウィンドウ左端の連続区間 $[s, e]$ で与えられるとき，
対応する変化点は開始時刻 $s$ と，存在すれば終了直後の $e+1$ である．
\end{definition}

\subsection{イベント帰属問題}

\begin{definition}[イベント帰属問題]\label{def:attribution}
頻出パターン集合 $\mathcal{F}$ とそのサポート時系列 $\{s_P\}_{P \in \mathcal{F}}$，
外部イベント集合 $\mathcal{E}$，有意水準 $\alpha$ が与えられたとき，
\emph{イベント帰属問題}とは，
$s_P$ の閾値交差変化点がイベント $e$ の時間窓と統計的に有意に関連する
全ての組 $(P, e) \in \mathcal{F} \times \mathcal{E}$ を，
偽発見率（FDR）を水準 $\alpha$ 以下に制御しながら同定することである．
\end{definition}

この問題の主要な課題は以下の3点である:

\begin{enumerate}
\item \textbf{規模}: $|\mathcal{F}|$ 個のパターンと $|\mathcal{E}|$ 個のイベントに対して
仮説数が $O(|\mathcal{F}| \times |\mathcal{E}|)$ となり，
厳格な多重検定補正が必要となる．

\item \textbf{副次パターン}: イベントがアイテム $i$ の頻度を増加させると，
$i$ を含む\emph{全て}のパターンのサポートが相関して変動し，偽陽性が大量に発生する．

\item \textbf{時間的混同}: 無関係なイベントの時間窓が重複する場合，帰属が曖昧になる．
\end{enumerate}
